\documentclass[12pt,a4paper]{article}

\usepackage[T1]{fontenc}
\usepackage[utf8]{inputenc}
\usepackage[french]{babel}
\usepackage{geometry}
\usepackage{graphicx}
\usepackage{hyperref}
\usepackage{tikz}
\usepackage[most]{tcolorbox}
\usepackage{listings}
\usepackage{xcolor}
\usepackage{array}
\usepackage{longtable}
\usepackage{fancyhdr}
\usepackage{titlesec}
\usepackage{enumitem}
\usepackage{setspace}
\usepackage{microtype}

\geometry{margin=2.4cm}
\setstretch{1.12}
\pagestyle{fancy}
\fancyhf{}

\definecolor{bg}{HTML}{F7FAFC}
\definecolor{panel}{HTML}{FFFFFF}
\definecolor{ink}{HTML}{0F172A}
\definecolor{muted}{HTML}{475569}
\definecolor{accent}{HTML}{00C2FF}
\definecolor{accent2}{HTML}{7C3AED}
\definecolor{accent3}{HTML}{14B8A6}
\definecolor{codebg}{HTML}{F1F5F9}
\definecolor{codeframe}{HTML}{CBD5E1}
\definecolor{line}{HTML}{0EA5E9}
\pagecolor{bg}
\color{ink}

\lhead{\textcolor{accent}{RobotLang}}
\rhead{\textcolor{muted}{Mini-projet Techniques de Compilation}}
\cfoot{\textcolor{muted}{\thepage}}

\hypersetup{
  colorlinks=true,
  linkcolor=accent2,
  urlcolor=accent,
  citecolor=accent3
}

\lstdefinestyle{robotstyle}{
  backgroundcolor=\color{codebg},
  frame=single,
  rulecolor=\color{codeframe},
  basicstyle=\ttfamily\small,
  keywordstyle=\color{accent2}\bfseries,
  commentstyle=\color{muted},
  stringstyle=\color{accent3},
  showstringspaces=false,
  breaklines=true,
  columns=fullflexible
}
\lstset{style=robotstyle}

\titleformat{\section}
  {\Large\bfseries\color{ink}}
  {\thesection}
  {0.8em}
  {\tikz[baseline=-0.6ex]\fill[accent] (0,0) rectangle (0.12,0.32);\hspace{0.5em}}
\titleformat{\subsection}
  {\large\bfseries\color{ink}}
  {\thesubsection}
  {0.8em}
  {}
\titlespacing*{\section}{0pt}{1.2em}{0.6em}
\titlespacing*{\subsection}{0pt}{1.0em}{0.45em}

\newtcolorbox{futuristicbox}[1][]{
  enhanced,
  colback=panel,
  colframe=accent,
  boxrule=0.8pt,
  arc=3mm,
  left=1.2mm,
  right=1.2mm,
  top=1.2mm,
  bottom=1.2mm,
  drop shadow={black!10!white},
  #1
}

\newtcolorbox{darkbox}[1][]{
  enhanced,
  colback=ink,
  colframe=accent3,
  coltext=white,
  boxrule=0.8pt,
  arc=3mm,
  left=1.2mm,
  right=1.2mm,
  top=1.2mm,
  bottom=1.2mm,
  drop shadow={black!18!white},
  #1
}

\newcommand{\smallcapslabel}[1]{\textcolor{accent}{\textsc{#1}}}

\title{\textbf{Rapport de mini-projet}\\\vspace{0.25cm}RobotLang : compilateur et interpr\`eteur pour un langage de robot virtuel}
\author{Nom de l'\`etudiant : \hspace{1cm} \textit{À compléter}}
\date{Année universitaire 2024--2025}

\begin{document}

\begin{titlepage}
  \thispagestyle{empty}
  \begin{tikzpicture}[remember picture,overlay]
    \fill[ink] (current page.north west) rectangle ([yshift=-7.1cm]current page.north east);
    \fill[accent] (current page.north west) rectangle ([yshift=-0.22cm]current page.north east);
    \fill[accent2,opacity=0.12] ([xshift=10.2cm,yshift=-1.6cm]current page.north west) circle (4.3cm);
    \fill[accent,opacity=0.08] ([xshift=15.6cm,yshift=-9.7cm]current page.north west) circle (5.2cm);
    \draw[accent,opacity=0.32,line width=1.2pt] ([xshift=1.2cm,yshift=-1.2cm]current page.north west) -- ([xshift=1.2cm,yshift=-13.1cm]current page.north west);
    \draw[accent3,opacity=0.32,line width=1.2pt] ([xshift=19.1cm,yshift=-1.2cm]current page.north west) -- ([xshift=19.1cm,yshift=-13.1cm]current page.north west);
  \end{tikzpicture}

  \vspace*{1.2cm}
  \begin{center}
    {\smallcapslabel{MODULE : TECHNIQUES DE COMPILATION}\par}
    \vspace{0.25cm}
    {\large\color{white}\textbf{DÉPARTEMENT GLSI - ISI}\par}
    \vspace{1.15cm}
    {\fontsize{42}{48}\selectfont\bfseries\color{accent}\textsf{RobotLang}\par}
    \vspace{0.35cm}
    {\Large\bfseries\color{white}\textsf{Conception et réalisation d'un compilateur pour un langage de robot virtuel}\par}
    \vspace{1.0cm}
    \begin{futuristicbox}[width=0.88\textwidth,colback=panel]
      \centering
      \begin{tabular}{>{\bfseries\color{ink}}p{5.3cm} p{7.0cm}}
        Sujet & Mini-projet de compilation / interprétation \\
        Technologie principale & Python \\
        Front-end & Analyse lexicale, syntaxique et sémantique \\
        Back-end & Bytecode et machine virtuelle \\
        Interface & CLI + GUI Tkinter \\
      \end{tabular}
    \end{futuristicbox}
    \vspace{0.8cm}
    \begin{darkbox}[width=0.88\textwidth]
      \centering
      {\large\textbf{Année universitaire 2024--2025}\par}
      \vspace{0.2cm}
      {\color{white!80}Un mini-compilateur visuel, modulaire et démonstratif\par}
    \end{darkbox}
  \end{center}
\end{titlepage}

\tableofcontents
\newpage

\section{Introduction}

\begin{futuristicbox}
Ce projet consiste à concevoir et réaliser \textbf{RobotLang}, un langage de programmation
dédié au pilotage d'un robot virtuel sur une grille 2D. L'objectif est de montrer une
chaîne complète de compilation, depuis le texte source jusqu'à l'exécution sur une machine
virtuelle.
\end{futuristicbox}

RobotLang répond au cahier des charges du mini-projet avec une approche plus visuelle et
plus originale qu'un langage abstrait classique. Le robot virtuel permet d'illustrer
facilement les structures de contrôle, les fonctions, les erreurs et l'exécution pas à pas.

L'idée générale est de fournir un mini-compilateur qui soit à la fois :
\begin{itemize}[nosep]
  \item clair à expliquer ;
  \item riche en fonctionnalités ;
  \item démontrable en direct ;
  \item suffisamment technique pour montrer une vraie architecture de compilation.
\end{itemize}

\section{Objectifs et périmètre}

Le projet vise à couvrir :
\begin{itemize}[nosep]
  \item l'analyse lexicale ;
  \item l'analyse syntaxique ;
  \item l'analyse sémantique ;
  \item l'optimisation d'expressions ;
  \item la compilation en bytecode ;
  \item l'exécution sur une machine virtuelle ;
  \item la visualisation graphique du résultat.
\end{itemize}

\begin{futuristicbox}
\smallcapslabel{Apport principal.} Le projet ne se limite pas à un interpréteur.
Il met en place une architecture découpée en modules, ce qui rend la solution plus proche
d'un vrai compilateur pédagogique.
\end{futuristicbox}

\section{Présentation du langage}

RobotLang est un DSL orienté robot. Le programmeur écrit des instructions comme
\texttt{move}, \texttt{turn}, \texttt{goto}, \texttt{randomize} ou \texttt{walkthrough}
dans une logique très proche des langages de script utilisés pour l'automatisation.

\subsection{Exemple de programme}

\begin{lstlisting}[caption={Exemple de programme RobotLang}]
let step = 1;
let nums = [1, 2, 3];

proc walk(n) {
  let i = 0;
  while i < n {
    move step;
    i = i + 1;
  }
  return n;
}

walk(3);
randomize;
goto(6, 6);
print at_goal;
\end{lstlisting}

\subsection{Fonctionnalités}

\begin{longtable}{|p{4cm}|p{10cm}|}
\hline
\textbf{Fonctionnalité} & \textbf{Description} \\
\hline
Variables & Déclaration avec \texttt{let} et affectation \\
\hline
Expressions & Arithmétique, comparaisons et logique \\
\hline
Contrôle & \texttt{if}, \texttt{else}, \texttt{while}, \texttt{repeat} \\
\hline
Procédures & Définition, appel et \texttt{return} \\
\hline
Listes & Littéraux de listes et indexation \\
\hline
Robot & \texttt{move}, \texttt{turn}, \texttt{wait}, \texttt{goto}, \texttt{pathfind} \\
\hline
Environnement & \texttt{randomize} pour reshuffle le monde \\
\hline
I/O & Lecture et écriture de fichiers \\
\hline
Interface & CLI, AST, tokens, GUI, walkthrough \\
\hline
\end{longtable}

\section{Architecture}

\begin{futuristicbox}
L'architecture suit le schéma classique d'un compilateur moderne :
\textbf{source text} $\rightarrow$ \textbf{tokens} $\rightarrow$ \textbf{AST}
$\rightarrow$ \textbf{analyse sémantique} $\rightarrow$ \textbf{bytecode}
$\rightarrow$ \textbf{machine virtuelle} $\rightarrow$ \textbf{effet visuel}.
\end{futuristicbox}

\subsection{Organisation du dépôt}

\begin{itemize}[nosep]
  \item \texttt{robotlang/tokenizer.py} : lexing ;
  \item \texttt{robotlang/parser.py} : parsing ;
  \item \texttt{robotlang/analyzer.py} : vérifications sémantiques ;
  \item \texttt{robotlang/optimizer.py} : constant folding ;
  \item \texttt{robotlang/compiler.py} : génération de bytecode ;
  \item \texttt{robotlang/vm.py} : exécution ;
  \item \texttt{robotlang/gui.py} : interface futuriste ;
  \item \texttt{tests/} : validation ;
  \item \texttt{examples/} : démonstration.
\end{itemize}

\subsection{Choix technologiques}

Python a été retenu pour sa simplicité, sa lisibilité et son intégration facile avec
Tkinter. Ce choix permet de développer rapidement une architecture lisible tout en restant
accessible pour un mini-projet de compilation.

\section{Analyse lexicale}

Le lexer lit le texte source caractère par caractère. Il produit des jetons avec position
source afin de fournir des erreurs précises.

\begin{itemize}[nosep]
  \item mots-clés : \texttt{let}, \texttt{if}, \texttt{while}, \texttt{proc}, \texttt{return}, etc. ;
  \item identifiants ;
  \item nombres ;
  \item chaînes ;
  \item symboles ;
  \item commentaires ;
  \item erreurs lexicales.
\end{itemize}

\begin{darkbox}
\smallcapslabel{Message d'erreur.} Le lexer signale les caractères inconnus avec ligne,
colonne et contexte visuel, ce qui rend la correction beaucoup plus simple.
\end{darkbox}

\section{Analyse syntaxique et AST}

Le parseur transforme les jetons en arbre syntaxique abstrait (AST). L'AST sépare la
syntaxe concrète du traitement sémantique et de la génération de bytecode.

\subsection{Nœuds principaux}

\begin{itemize}[nosep]
  \item \texttt{Program}, \texttt{Block} ;
  \item \texttt{LetStmt}, \texttt{AssignStmt}, \texttt{SetIndexStmt} ;
  \item \texttt{IfStmt}, \texttt{WhileStmt}, \texttt{RepeatStmt} ;
  \item \texttt{ProcDef}, \texttt{ReturnStmt}, \texttt{BreakStmt}, \texttt{ContinueStmt} ;
  \item \texttt{MoveStmt}, \texttt{TurnStmt}, \texttt{GotoStmt}, \texttt{RandomizeStmt} ;
  \item \texttt{BinaryOp}, \texttt{UnaryOp}, \texttt{ListLiteral}, \texttt{IndexExpr} ;
  \item \texttt{CallExpr}, \texttt{ReadExpr}, \texttt{PathFindExpr}.
\end{itemize}

\subsection{Grammaire simplifiée}

\begin{lstlisting}[caption={Grammaire simplifiée}]
program     -> statement* EOF
statement   -> letStmt | assignStmt | ifStmt | whileStmt | repeatStmt
            | procDef | returnStmt | breakStmt | continueStmt
            | moveStmt | turnStmt | gotoStmt | randomizeStmt
            | printStmt | writeStmt | exprStmt

expression  -> logicOr
logicOr     -> logicAnd ("or" logicAnd)*
logicAnd    -> equality ("and" equality)*
equality    -> comparison (("==" | "!=") comparison)*
comparison  -> term (("<" | "<=" | ">" | ">=") term)*
term        -> factor (("+" | "-") factor)*
factor      -> unary (("*" | "/" | "%") unary)*
unary       -> ("-" | "!" | "not") unary | postfix
postfix     -> primary ("(" args? ")" | "[" expression "]")*
primary     -> number | string | bool | ident | list | predicate
\end{lstlisting}

\section{Analyse sémantique}

L'analyse sémantique vérifie la cohérence du programme :
\begin{itemize}[nosep]
  \item variables déclarées avant usage ;
  \item procédures déclarées avant appel ;
  \item nombre d'arguments correct ;
  \item \texttt{return} uniquement dans une procédure ;
  \item \texttt{break} et \texttt{continue} uniquement dans une boucle ;
  \item prédicats intégrés valides ;
  \item cohérence des indices et des types.
\end{itemize}

\begin{futuristicbox}
Les portées imbriquées sont gérées par des scopes hiérarchiques. Cela permet de garder
une séparation claire entre variables locales, paramètres et contexte global.
\end{futuristicbox}

\section{Optimisation}

Une passe de \textbf{constant folding} simplifie les expressions entièrement constantes.

\begin{lstlisting}[language=Python,caption={Constant folding}]
let x = 2 + 3 * 4;
\end{lstlisting}

peut être transformé avant exécution en :

\begin{lstlisting}[language=Python]
let x = 14;
\end{lstlisting}

\section{Bytecode et VM}

Le projet compile l'AST en bytecode puis exécute ce bytecode dans une machine virtuelle.
Cette séparation clarifie le pipeline et permet d'avoir un vrai back-end d'exécution.

\subsection{Instructions}

\begin{itemize}[nosep]
  \item \texttt{PUSH\_CONST}, \texttt{LOAD}, \texttt{STORE}, \texttt{DECLARE} ;
  \item \texttt{CALL}, \texttt{RETURN} ;
  \item \texttt{BINARY}, \texttt{UNARY} ;
  \item \texttt{MOVE}, \texttt{TURN}, \texttt{WAIT}, \texttt{GOTO}, \texttt{RANDOMIZE} ;
  \item \texttt{JUMP}, \texttt{JUMP\_IF\_FALSE}, \texttt{JUMP\_IF\_TRUE} ;
  \item \texttt{ENTER\_SCOPE}, \texttt{EXIT\_SCOPE}.
\end{itemize}

\subsection{Intérêt pédagogique}

\begin{darkbox}
\smallcapslabel{Pourquoi le bytecode ?} Il montre très bien l'idée de transformation
intermédiaire avant exécution, ce qui correspond bien à un mini-projet de compilation.
\end{darkbox}

\section{Interface graphique}

\begin{futuristicbox}
L'interface Tkinter a été conçue comme un véritable \textbf{panneau de mission}.
Elle combine un éditeur de code, une grille visuelle, un journal d'exécution et une zone
d'explication en mode walkthrough.
\end{futuristicbox}

\subsection{Mode walkthrough}

Le bouton \texttt{WALKTHROUGH} exécute le pipeline avec des pauses visibles :
\begin{enumerate}[nosep]
  \item tokenisation ;
  \item parsing ;
  \item optimisation ;
  \item compilation en bytecode ;
  \item exécution pas à pas ;
  \item explication de chaque instruction.
\end{enumerate}

\subsection{Captures d'écran à insérer}

Remplacez les encadrés ci-dessous par vos propres captures de l'interface.

\begin{figure}[h!]
  \centering
  \begin{futuristicbox}[width=0.9\textwidth,height=5.2cm]
    \centering \vspace*{1.5cm}
    \textbf{Capture d'écran de l'interface principale}
  \end{futuristicbox}
  \caption{Interface principale de RobotLang}
\end{figure}

\begin{figure}[h!]
  \centering
  \begin{futuristicbox}[width=0.9\textwidth,height=5.2cm]
    \centering \vspace*{1.5cm}
    \textbf{Capture d'écran du mode walkthrough}
  \end{futuristicbox}
  \caption{Mode walkthrough}
\end{figure}

\begin{figure}[h!]
  \centering
  \begin{futuristicbox}[width=0.9\textwidth,height=5.2cm]
    \centering \vspace*{1.5cm}
    \textbf{Capture d'écran d'un programme en exécution}
  \end{futuristicbox}
  \caption{Exécution d'un programme RobotLang}
\end{figure}

\section{Monde virtuel}

Le robot évolue sur une grille 2D contenant :
\begin{itemize}[nosep]
  \item une position de départ ;
  \item une direction ;
  \item des obstacles ;
  \item un objectif ;
  \item des règles de collision et de bornes.
\end{itemize}

\subsection{Commandes spécifiques}

\begin{itemize}[nosep]
  \item \texttt{move} : avance ;
  \item \texttt{turn} : rotation par multiples de 90 degrés ;
  \item \texttt{goto(x, y)} : déplacement planifié ;
  \item \texttt{pathfind(x, y)} : calcul de chemin ;
  \item \texttt{randomize} : nouvelle mission aléatoire.
\end{itemize}

\section{Tests et validation}

Les tests unitaires couvrent :
\begin{itemize}[nosep]
  \item le programme de démonstration ;
  \item les variables et affectations ;
  \item \texttt{while}, \texttt{break}, \texttt{continue} ;
  \item \texttt{randomize} ;
  \item la planification de chemin ;
  \item les erreurs lexicales, syntaxiques, sémantiques et runtime.
\end{itemize}

\begin{futuristicbox}
Les commandes de validation utilisées sont :
\begin{lstlisting}[language=bash]
py -B -m unittest discover -s tests
py -B -m robotlang --demo
py -B -m robotlang --ast examples/demo.rbt
py -B -m robotlang --gui
\end{lstlisting}
\end{futuristicbox}

\section{Difficultés et solutions}

\begin{itemize}[nosep]
  \item \textbf{Erreurs lisibles} : stockage de la ligne et de la colonne pour chaque jeton.
  \item \textbf{Portées imbriquées} : frames hiérarchiques dans l'analyse et la VM.
  \item \textbf{Walkthrough} : événements d'exécution et pauses contrôlées.
  \item \textbf{Design} : séparation claire entre le calcul et la présentation.
\end{itemize}

\section{Bilan}

\begin{darkbox}
RobotLang dépasse le minimum demandé par le cahier des charges : il comporte une vraie
chaîne de traitement, une optimisation, une machine virtuelle, une interface graphique et
des fonctions bonus originales.
\end{darkbox}

\section{Conclusion}

Ce mini-projet illustre les principales étapes d'un compilateur : analyse lexicale,
analyse syntaxique, vérification sémantique, optimisation, compilation et exécution.
Le thème du robot virtuel donne au projet une identité visuelle forte et un côté
présentable en démonstration.

RobotLang constitue une base solide pour aller plus loin : ajout de coloration syntaxique,
génération de code natif, animations plus fluides ou encore éditeur intégré.

\appendix
\section{Annexe : commandes principales}

\begin{itemize}[nosep]
  \item \texttt{let x = 5;}
  \item \texttt{if obstacle\_ahead \{ ... \}}
  \item \texttt{while i < 10 \{ ... \}}
  \item \texttt{proc nom(param) \{ ... \}}
  \item \texttt{return valeur;}
  \item \texttt{randomize;}
  \item \texttt{goto(6, 6);}
\end{itemize}

\section{Annexe : programme de démonstration}

\begin{lstlisting}
let step = 1;
let path = [1, 2, 3];

proc turn_around() {
  turn 90;
  turn 90;
}

proc walk_square(n) {
  let i = 0;
  while i < n {
    move step;
    i = i + 1;
  }
  return n;
}

walk_square(2);
turn_around();
if obstacle_ahead {
  goto(0, 1);
}
goto(6, 6);
print at_goal;
\end{lstlisting}

\section{Annexe : insertion de captures}

Créez un dossier \texttt{screenshots/} puis remplacez les blocs de la section GUI par :

\begin{lstlisting}[language=TeX]
\includegraphics[width=0.9\textwidth]{screenshots/gui_main.png}
\includegraphics[width=0.9\textwidth]{screenshots/gui_walkthrough.png}
\includegraphics[width=0.9\textwidth]{screenshots/gui_run.png}
\end{lstlisting}

\end{document}
